\chapter{Accesso ai dati da software}

\section{PL/pgSQL e Stored Procedure}
\textbf{PL/pgSQL} è un linguaggio di programmazione procedurale
supportato da PostgreSQL, che permette di eseguire operazioni
complesse direttamente all'interno del DBMS.\\
A differenza di SQL, che è un linguaggio dichiarativo
(dichiariamo \textit{cosa} vogliamo ottenere), PL/pgSQL
permette di specificare anche \textit{come} eseguire
le operazioni, utilizzando:
\begin{itemize}
    \item Variabili locali;
    \item Cicli (\texttt{FOR}, \texttt{WHILE});
    \item Strutture condizionali (\texttt{IF}, \texttt{ELSE});
    \item Query SQL all'interno del codice.
\end{itemize}
\noindent Queste funzionalità permettono di implementare logiche
complesse che sarebbero difficili da esprimere
attraverso semplici query SQL.
\subsection{Funzioni e Stored Procedure}
In PostgreSQL possiamo definire delle \textbf{funzioni},
ovvero blocchi di codice memorizzati nel database
ed eseguibili ogni volta che vengono richiamati.\\
Una funzione può:
\begin{itemize}
    \item Ricevere parametri in ingresso;
    \item Eseguire operazioni sui dati;
    \item Restituire un valore oppure una tabella;
    \item Essere richiamata tramite istruzioni SQL;
    \item Essere utilizzata da un trigger.
\end{itemize}
\noindent Una volta creata, la funzione viene memorizzata
nel database e può essere riutilizzata senza
dover riscrivere ogni volta le istruzioni.\\
Una funzione con \texttt{RETURNS VOID} non restituisce
un valore utile, ma rimane comunque una funzione.
\subsection{Struttura di una funzione PL/pgSQL}
La sintassi generale è:
\begin{verbatim}
CREATE OR REPLACE FUNCTION nome_funzione(arg1 arg_type, .., argn arg_type)
RETURNS {return_type | TABLE(clo1 col_type,...) | TRIGGER} AS
$$
DECLARE
    -- Dichiarazione delle variabili
BEGIN
    -- Istruzioni da eseguire
END;
$$ LANGUAGE plpgsql;
\end{verbatim}

Analizziamo le parti principali:
\begin{itemize}
    \item \texttt{CREATE FUNCTION}: crea una nuova funzione.
    \item \texttt{OR REPLACE}: permette di sostituire
    una funzione già esistente con la stessa
    identità (nome e tipi dei parametri di ingresso).
    \item \texttt{parametri}: valori ricevuti
    in ingresso dalla funzione.
    \item \texttt{RETURNS}: specifica il tipo
    di risultato restituito.
    \item \texttt{\$\$}: delimita il corpo della funzione.
    \item \texttt{DECLARE}: sezione opzionale
    in cui vengono dichiarate le variabili locali.
    \item \texttt{BEGIN ... END}: contiene
    le istruzioni da eseguire.
    \item \texttt{LANGUAGE plpgsql}: specifica
    il linguaggio utilizzato.
\end{itemize}

Il tipo di ritorno può essere:
\begin{itemize}
    \item \texttt{RETURNS INTEGER}: restituisce un intero.
    \item \texttt{RETURNS REAL}: restituisce un numero reale.
    \item \texttt{RETURNS VOID}: non restituisce
    un valore utile.
    \item \texttt{RETURNS TABLE(...)}: restituisce
    una tabella con colonne specificate.
    \item \texttt{RETURNS TRIGGER}: definisce
    una funzione utilizzabile da un trigger.
\end{itemize}
\textcolor{verdeEsempio}{Esempio:} calcolo del totale delle fatture
Consideriamo la relazione:

\begin{verbatim}
Fattura(
    soggetto VARCHAR(100),
    imponibile REAL,
    aliquota INTEGER
)
\end{verbatim}
\noindent 
dove:
\begin{itemize}
    \item \texttt{soggetto}: identifica il soggetto
    a cui è associata la fattura.
    \item \texttt{imponibile}: importo della fattura
    prima dell'applicazione dell'imposta.
    \item \texttt{aliquota}: percentuale di imposta
    applicata.
\end{itemize}
\noindent
Supponiamo che la tabella contenga:

\begin{center}
\begin{tabular}{|c|c|c|}
\hline
Soggetto & Imponibile & Aliquota \\
\hline
a & 1000 & 22 \\
a & 30 & 4 \\
b & 10000 & 10 \\
\hline
\end{tabular}
\end{center}
\noindent
Vogliamo creare una funzione che, dato un soggetto,
restituisca una tabella contenente il suo identificativo
e il totale delle fatture comprensivo delle imposte.\\
Per ogni fattura il totale viene calcolato come:

\[
\text{Totale fattura} =
\text{Imponibile}\cdot
\frac{100+\text{Aliquota}}{100}
\]

Il totale complessivo è quindi:

\[
T(s)=
\sum_{\substack{f\in Fattura\\f.soggetto=s}}
f.imponibile\cdot
\frac{100+f.aliquota}{100}
\]

\textbf{Implementazione in PL/pgSQL}

\begin{verbatim}
CREATE OR REPLACE FUNCTION totale_fattura(sogg VARCHAR)
RETURNS TABLE (sgt VARCHAR, tot REAL) AS
$$
DECLARE
    r RECORD;
    totale REAL;
BEGIN
    totale := 0;

    FOR r IN
        SELECT *
        FROM Fattura
        WHERE soggetto = sogg
    LOOP
        totale := totale +
            r.imponibile * (100 + r.aliquota) / 100;
    END LOOP;

    RETURN QUERY SELECT sogg, totale;
END;
$$ LANGUAGE plpgsql;
\end{verbatim}
\textbf{Spiegazione passo per passo:}

\begin{enumerate}
    \item \texttt{totale\_fattura(sogg VARCHAR)}:\\
    La funzione riceve il nome del soggetto
    di cui vogliamo calcolare il totale.

    \item \texttt{RETURNS TABLE (sgt VARCHAR, tot REAL)}:\\
    La funzione restituisce una tabella con due colonne:
    il soggetto e il totale calcolato.

    \item \texttt{r RECORD}:\\
    Dichiara una variabile capace di contenere
    una riga restituita da una query.

    \item \texttt{totale REAL}:\\
    Dichiara una variabile utilizzata
    per accumulare il totale.

    \item \texttt{totale := 0}:\\
    Inizializza il totale a zero.
    L'operatore \texttt{:=} rappresenta
    l'assegnazione in PL/pgSQL.

    \item \texttt{FOR r IN SELECT ... LOOP}:\\
    Scorre tutte le fatture appartenenti
    al soggetto passato come parametro.\\
    A ogni iterazione, \texttt{r} contiene
    una diversa riga della query.

    \item \texttt{totale := totale + ...}:\\
    Calcola l'importo comprensivo di imposta
    della fattura corrente e lo aggiunge
    al totale accumulato.

    \item \texttt{END LOOP}:\\
    Termina il ciclo dopo aver elaborato
    tutte le fatture del soggetto.

    \item \texttt{RETURN QUERY SELECT sogg, totale}:\\
    Restituisce una riga contenente
    il soggetto e il totale calcolato.
\end{enumerate}
\noindent \textbf{Esecuzione della funzione}
Per richiamare la funzione scriviamo:
\begin{verbatim}
SELECT * FROM totale_fattura('a');
\end{verbatim}
Per il soggetto \texttt{a} abbiamo due fatture:

\[
1000\cdot\frac{100+22}{100}=1220
\]

\[
30\cdot\frac{100+4}{100}=31.2
\]

Quindi:

\[
T(a)=1220+31.2=\boxed{1251.2}
\]

Il risultato sarà:

\begin{center}
\begin{tabular}{|c|c|}
\hline
sgt & tot \\
\hline
a & 1251.2 \\
\hline
\end{tabular}
\end{center}

\subsection{Alternativa: calcolo tramite SQL}

L'operazione precedente può essere eseguita
anche direttamente attraverso una query SQL:

\begin{verbatim}
SELECT soggetto,
       SUM(imponibile * (100 + aliquota) / 100)
           AS totale
FROM Fattura
WHERE soggetto = 'a'
GROUP BY soggetto;
\end{verbatim}
Questa query produce lo stesso risultato
senza utilizzare cicli o variabili.\\
Infatti, quando un'operazione può essere espressa
facilmente attraverso SQL, generalmente è preferibile
utilizzare una query rispetto a un ciclo procedurale.

\subsection{Vantaggi delle Stored Procedure}

Anche se molte operazioni possono essere realizzate
direttamente tramite SQL, le funzioni e le procedure
memorizzate nel database offrono importanti vantaggi.

\begin{enumerate}
    \item \textbf{Riutilizzabilità del codice:}\\
    La logica viene definita una sola volta
    e può essere richiamata più volte.

    \item \textbf{Maggiore sicurezza:}\\
    È possibile concedere agli utenti
    il permesso di eseguire determinate funzioni
    senza concedere necessariamente l'accesso
    diretto alle tabelle sottostanti.\\
    Questo richiede una corretta configurazione
    dei privilegi e delle funzioni.

    \item \textbf{Coerenza dei dati:}\\
    Le operazioni vengono eseguite secondo
    regole prestabilite, riducendo il rischio
    che gli utenti eseguano modifiche
    non coerenti con la logica applicativa.

    \item \textbf{Centralizzazione della logica:}\\
    Le regole applicative possono essere
    implementate direttamente nel DBMS,
    evitando di duplicarle nei diversi programmi
    che accedono al database.

    \item \textbf{Integrazione con i trigger:}\\
    Le funzioni possono essere associate
    a trigger per eseguire automaticamente
    determinate operazioni quando si verificano
    eventi sul database.
\end{enumerate}
\noindent \textbf{Concetto fondamentale per l'esame:}\\
SQL è un linguaggio dichiarativo, mentre PL/pgSQL
aggiunge funzionalità procedurali come variabili,
cicli e condizioni.\\
Le funzioni memorizzate nel DBMS permettono
di incapsulare e riutilizzare operazioni,
centralizzare la logica applicativa e,
attraverso un'adeguata gestione dei privilegi,
migliorare la sicurezza e la coerenza dei dati.
\subsection{Trigger in PL/pgSQL}

Un \textbf{trigger} è un meccanismo del DBMS che permette
di eseguire automaticamente una funzione quando si verifica
un determinato evento su una tabella.\\
Gli eventi principali che possono attivare un trigger sono:
\begin{itemize}
    \item \texttt{INSERT}: inserimento di nuove tuple;
    \item \texttt{UPDATE}: modifica di tuple esistenti;
    \item \texttt{DELETE}: eliminazione di tuple.
\end{itemize}
\noindent
Un trigger viene quindi \textbf{associato a una tabella}
e si attiva automaticamente quando si verifica
l'evento per cui è stato definito.\\
I trigger sono particolarmente utili per:
\begin{itemize}
    \item Controllare la validità dei dati;
    \item Imporre vincoli e regole applicative;
    \item Aggiornare automaticamente altre informazioni;
    \item Registrare le modifiche effettuate nel database.
\end{itemize}
\noindent
A differenza di una normale funzione, che viene richiamata
esplicitamente tramite SQL, un trigger viene eseguito
automaticamente dal DBMS.

\subsubsection{Sintassi di un trigger}

La sintassi generale in PostgreSQL è:

\begin{verbatim}
CREATE TRIGGER nome_trigger [ BEFORE | AFTER | INSTEAD OF ]
[ INSERT | UPDATE | DELETE ] ON nome_tabella
[ FOR EACH ROW | FOR EACH STATEMENT ]
EXECUTE FUNCTION nome_funzione();
\end{verbatim}

\noindent
La funzione richiamata dal trigger deve essere stata
precedentemente definita con \texttt{RETURNS TRIGGER}.\\
La creazione di un trigger richiede quindi due passaggi:
\begin{enumerate}
    \item Definire la funzione che contiene
    le operazioni da eseguire;
    \item Creare il trigger che associa tale funzione
    a un evento su una tabella.
\end{enumerate}

\subsubsection{Semantica dei trigger}

La definizione di un trigger specifica tre aspetti
fondamentali:
\begin{enumerate}
    \item \textbf{Quando} deve essere eseguito;
    \item \textbf{Quale evento} ne provoca l'attivazione;
    \item \textbf{Quante volte} deve essere eseguito.
\end{enumerate}

\noindent
\textbf{1. Momento di esecuzione: BEFORE, AFTER, INSTEAD OF}

\begin{itemize}
    \item \texttt{BEFORE}: il trigger viene eseguito
    prima che l'operazione di modifica venga effettuata.\\
    È utile per controllare o modificare i dati
    prima che vengano memorizzati.

    \item \texttt{AFTER}: il trigger viene eseguito
    dopo che l'operazione di modifica è stata effettuata.\\
    È utile, ad esempio, per registrare modifiche
    o aggiornare informazioni correlate.

    \item \texttt{INSTEAD OF}: il trigger viene eseguito
    al posto dell'operazione richiesta.\\
    In PostgreSQL viene utilizzato sulle viste
    e permette di definire come gestire
    operazioni di modifica su di esse.
\end{itemize}

\noindent
\textbf{2. Evento di attivazione: INSERT, UPDATE, DELETE}
Il trigger può essere attivato quando viene eseguita
una delle seguenti operazioni:
\begin{itemize}
    \item \texttt{INSERT}: inserimento di una tupla;
    \item \texttt{UPDATE}: aggiornamento di una tupla;
    \item \texttt{DELETE}: eliminazione di una tupla.
\end{itemize}
\noindent
È possibile associare lo stesso trigger
anche a più eventi.\medskip\\
\textbf{3. Frequenza di esecuzione: ROW e STATEMENT}

\begin{itemize}
    \item \texttt{FOR EACH ROW}: il trigger viene eseguito
    una volta per ogni tupla coinvolta nell'operazione.

    \item \texttt{FOR EACH STATEMENT}: il trigger viene
    eseguito una sola volta per ogni istruzione SQL,
    indipendentemente dal numero di tuple coinvolte.\\
    È il comportamento predefinito.
\end{itemize}

\noindent
\textbf{Esempio:}\\
Supponiamo di eseguire:

\begin{verbatim}
UPDATE Studenti
SET voto = voto + 1
WHERE corso = 'Basi di Dati';
\end{verbatim}

Se l'istruzione modifica 10 tuple:
\begin{itemize}
    \item Con \texttt{FOR EACH ROW}, la funzione
    del trigger viene eseguita \textbf{10 volte},
    una per ogni tupla modificata.

    \item Con \texttt{FOR EACH STATEMENT}, la funzione
    del trigger viene eseguita \textbf{una sola volta},
    indipendentemente dalle 10 tuple modificate.
\end{itemize}
\noindent
\textbf{Attenzione:} un trigger
\texttt{FOR EACH STATEMENT} viene eseguito anche
se l'istruzione non modifica alcuna tupla.\\
Un trigger \texttt{FOR EACH ROW}, invece,
non viene eseguito se nessuna tupla è coinvolta.

\subsubsection{Variabili speciali NEW e OLD}

All'interno delle funzioni associate ai trigger
\texttt{FOR EACH ROW}, PostgreSQL mette a disposizione
due variabili speciali di tipo \texttt{RECORD}:

\begin{itemize}
    \item \texttt{NEW}: rappresenta la \textbf{nuova tupla},
    cioè i valori dopo un inserimento o una modifica.

    \item \texttt{OLD}: rappresenta la \textbf{vecchia tupla},
    cioè i valori prima di una modifica o eliminazione.
\end{itemize}

\noindent
La disponibilità delle variabili dipende
dall'operazione eseguita:

\begin{center}
\begin{tabular}{|c|c|c|}
\hline
\textbf{Operazione} & \textbf{NEW} & \textbf{OLD} \\
\hline
INSERT & Disponibile & NULL \\
\hline
UPDATE & Disponibile & Disponibile \\
\hline
DELETE & NULL & Disponibile \\
\hline
\end{tabular}
\end{center}

\noindent
Nei trigger \texttt{FOR EACH STATEMENT},
entrambe le variabili sono \texttt{NULL},
poiché il trigger non fa riferimento
a una singola tupla.
\medskip\\
\textbf{Esempio con UPDATE}
Consideriamo la tabella:
\begin{verbatim}
Studenti(matricola, nome, voto)
\end{verbatim}
Supponiamo che una tupla contenga:

\[
(12345,\text{Mario},25)
\]
ed eseguiamo:

\begin{verbatim}
UPDATE Studenti
SET voto = 28
WHERE matricola = 12345;
\end{verbatim}
All'interno di un trigger \texttt{FOR EACH ROW}
associato a questo aggiornamento avremo:

\begin{itemize}
    \item \texttt{OLD.voto = 25}: valore precedente;
    \item \texttt{NEW.voto = 28}: nuovo valore.
\end{itemize}

\noindent
Quindi, grazie a \texttt{OLD} e \texttt{NEW},
possiamo confrontare i valori prima e dopo
una modifica.

\subsubsection{(da vedere se necessario in caso rimuovere) Esempio pratico: controllo dei voti}

Supponiamo di voler impedire l'inserimento
di voti non validi nella tabella \texttt{Studenti}.\\
Un voto deve essere compreso tra 18 e 30.\\
Possiamo definire una funzione trigger:

\begin{verbatim}
CREATE OR REPLACE FUNCTION controlla_voto()
RETURNS TRIGGER AS
$$
BEGIN
    IF NEW.voto < 18 OR NEW.voto > 30 THEN
        RAISE EXCEPTION 'Voto non valido';
    END IF;

    RETURN NEW;
END;
$$ LANGUAGE plpgsql;
\end{verbatim}

\noindent
Successivamente creiamo il trigger:

\begin{verbatim}
CREATE TRIGGER verifica_voto
BEFORE INSERT OR UPDATE
ON Studenti
FOR EACH ROW
EXECUTE FUNCTION controlla_voto();
\end{verbatim}

\noindent
\textbf{Funzionamento:}
\begin{enumerate}
    \item Viene eseguita un'istruzione
    \texttt{INSERT} oppure \texttt{UPDATE}
    sulla tabella \texttt{Studenti}.

    \item Prima della modifica, il trigger
    richiama automaticamente
    la funzione \texttt{controlla\_voto()}.

    \item La funzione controlla il valore
    di \texttt{NEW.voto}.

    \item Se il voto non è valido,
    \texttt{RAISE EXCEPTION} genera un errore
    e impedisce l'operazione.

    \item Se il voto è valido,
    \texttt{RETURN NEW} permette di proseguire
    con l'inserimento o l'aggiornamento.
\end{enumerate}

\noindent
\textbf{Importante:} nei trigger
\texttt{BEFORE FOR EACH ROW},
il valore restituito dalla funzione
influenza l'operazione:
\begin{itemize}
    \item \texttt{RETURN NEW}: permette di proseguire
    utilizzando la nuova tupla, eventualmente modificata.

    \item \texttt{RETURN NULL}: impedisce l'elaborazione
    della tupla corrente, senza necessariamente
    generare un errore.
\end{itemize}
\noindent
Nei trigger \texttt{AFTER}, invece,
il valore restituito dalla funzione viene ignorato.
\noindent
\textbf{Distinzione fondamentale:}\\
Una funzione normale viene richiamata
esplicitamente, mentre una funzione associata
a un trigger viene eseguita automaticamente
quando si verifica l'evento specificato.\\

SLIDE: fino 17